\chapter{Background}\label{ch:background}

\section{State of the Art Workflow}

\begin{wrapfigure}[9]{r}{0.48\linewidth}
    \vspace{-10pt}
    \includegraphics[width=\linewidth]{3dmodel_box.png}
    \caption{3D model of simple box}
    \label{simple-3d-model-box}
\end{wrapfigure}

When analyzing the process from idea to print, several main steps can be identified as shown in \ref{workflow}. 
Using a simple box with a snap-fit closure and hinge as an example, the following analysis will focus on the steps from exporting a model to printing, highlighting in particular how the workflow differs depending on the file format used in regards to three points.

\begin{enumerate}
    \item A model consisting of multiple parts 
    \item Using multiple materials (e.g. printing the hinge in a flexible material)
    \item Using model parts as modifiers (e.g. the light blue boxes as modifiers to print more wall loops in this area)
\end{enumerate} 

Exporting the model from Fusion is possible via two dialogs: the general export and the 3D print specific one.
The two dialogs each offer a different set of file formats due to their different purposes.
Since this analysis deals with 3D printing in particular, only the export dialog specific to 3D printing will be considered, which offers to export as STL, 3MF and OBJ format. 
For all three the user can decide whether to generate one file containing all selected components and bodies, or one file per body. 

\subsection{Exporting as STL}

Exporting the model as an STL file immediately shows one important distinction. 
``STL native files describe only the surface geometry of a three dimensional object without any representation of color, texture or other common CAD model attributes,''\cite{iancu_about_2018} meaning that even if a user applies materials in Fusion, that information won't be exported, and the user has to apply all materials anew in OrcaSlicer. 

Furthermore there will always be one body in the resulting file, even if the model is comprised of multiple bodies within Fusion. 
Therefore the user has to choose whether to combine all bodies into one or exporting each body as a separate file.
The first option results in one extra step to split the model into separate objects within OrcaSlicer to be able to print the hinges in a flexible material or to use parts as modifiers (e.g. the blue boxes for additional wall loops). 
Alternatively exporting each body as a separate file might loose their spatial relations to each other while keeping the bodies separate.
OrcaSlicer does ask the user whether the files should be loaded as a single object with multiple parts. 
The advantage over loading each file as an object is the preservation of the spatial relation and that parts, like the hinges in the example box, can be placed at any height while a single object is always placed with at least one surface flush with the printing plate. 

\subsection{Exporting as OBJ}

The OBJ file format ``is the preferred 3D printer file format for multicolor printing''.\cite{iancu_about_2018}
Consequently, OrcaSlicer shows a particular dialog (fig. \ref{fig:obj_color_dialog}) when opening an OBJ file, asking the user how they want to map the file's colors to current or new filament presets.
This makes it easy for users to align filament colors (i.e. the color the filament is represented as in OrcaSlicer, not necessarily the color of the filament itself) to colors applied in Fusion.
Just like STL files though, the model is generated as one body per file resulting in the same dialogs and extra steps.

\subsection{Exporting as 3MF}

Unlike the other two formats, 3MF not only preserves separate bodies even when exported as a single file making it more user-friendly, it also includes ``information about the model’s color, material, and textures.''\cite{iancu_about_2018}
Initially when importing a 3mf file containing multiple objects, OrcaSlicer asks the user whether the file should be loaded as a single object with multiple parts instead of considering them as multiple objects similarly to the dialog when importing multiple STL or OBJ files at once.

While Fusion does export materials as RGBA colors in one hexcode per body, OrcaSlicer does not recognize them so far. 
Consequently users have to reapply materials manually which has been criticized by users in OrcaSlicer's Issues.\cite{sheminasalam_single_nodate}

\subsection{Applying Settings with and without Presets}

Generally, materials can be set using system or user defined presets per object. 
Material settings entail for example the filament itself, nozzle temperature, general printing speed and cooling settings.
OrcaSlicer actively asks the user, whether they want to save any changes to material parameters they have made, which makes it easy to save either parameters specific to one's printer and filament or even specific to one project while simultaneously keeping default material presets untouched.

Furthermore settings like the amount of wall loops, infill percentage and pattern, and printing speed for single layers might need to be customized. 
It is possible to save these as presets as well, however only globally, \ie they can only be applied to the entire project, not to specific objects within.
As criticized in the \cite[OrcaSlicer Issue \#7769]{headamame_copypaste_2024}, particularly for projects containing multiple objects, using presets is not a valid option.


\section{Foundations}
\begin{itemize}
    \item what is CAD software?
    \item what is a slicer?

    ``Processing the 3D model occurred in specialized software known as a slicer (as it slices the model into layers for the printer), which allowed users to select print settings, such as layer height and print speed''\cite{hudson_understanding_2016}.
    \item 3MF File Format $\rightarrow$ why 3MF and not STL or other file formats? why not OBJ, since it supports exporting materials from fusion and mapping then in OrcaSlicer natively?
\end{itemize}

\subsubsection{The 3MF specification} \cite{3mf_consortium_3mf_2025}
\begin{itemize}
    \item ``3MF Document markup has been designed in anticipation of the evolution of this specification. It also allows third parties to extend the markup. Extensions are a critical part of 3MF, and as such, this core specification is as narrow as possible.'' [p.~8]
    \item 3MF = ZIP archive following the .ZIP File Format Specifications by PKWARE Inc. containing one or more 3D payloads
    \item 3D payload = ``complete collection of interdependent parts and relationships within a package'' consisting of ``a 3D object definition and its supporting files''
    \item 3MF document parts $\rightarrow$ table \ref{doc_parts}, figure \ref{doc_structure}

    \begin{table}
        \renewcommand{\arraystretch}{1.3}
        \begin{tabular}{m{0.21\linewidth} m{0.5\linewidth} m{0.14\linewidth}}
            \textbf{Name} & \textbf{Description} & \textbf{Required /} \textbf{Optional} \\
            \hline
            3D Model & Contains the description of one or more 3D objects for
manufacturing. & Required \\
            Core Properties & The OPC part that contains various document properties. & Optional \\
            Object Thumbnail & Contains a small JPEG or PNG image that represents a 3D object in a 3D Model. & Optional \\
            3D Texture & Contains a texture used to apply color to a 3D object in the 3D Model part (available for extensions) & Optional \\
            PrintTicket & Provides settings to be used when outputting the 3D object(s) in the 3D Model part. & Optional \\
            Custom Parts & OPC parts that are associated with metadata & Optional
        \end{tabular}\smallskip

        \caption{Excerpt of 3MF document parts\cite[p.~4]{3mf_consortium_3mf_2025}}
        \label{doc_parts}
    \end{table}
    
    \begin{figure}[htb]
    \dirtree{%
        .1 example.3mf.
        .2 \detokenize{_rels}/.
        .2 \detokenize{[Content_Types].xml}/.
        .2 3D/.
        .3 3dmodel.model.
        .2 Metadata/.
        .3 thumbnail.png.
    }
    \caption{Typical internal structure of 3MF file as exported from \emph{Autodesk Fusion}}
    \label{doc_structure}
\end{figure}

    \item PrintTicket $\rightarrow$ seems to be OS specific, so probably not helpful but \todo investigate further
    \item most important: 3D Model Part = root of 3D payload $\rightarrow$ ``contains definitions of one or more objects to be fabricated by 3D manufacturing processes''\cite[p.~6]{3mf_consortium_3mf_2025} $\rightarrow$ contains resource definitions (objects, materials) and items to be build 

    \begin{itemize}
        \item there MUST be exactly one <model> element in a 3D Model part
        \item a model must have two additional child elements: <resources> and <build>
        \item <resources> element
        
        provides a set of definitions that can be drawn from to define a 3D object; root element of a library of constituent pieces of the overall 3D object definition $\rightarrow$ objects, properties, materials; `object resource': 3d object that could be but will not be manufactured on its own
        
        \item <build> element
        
        provides a set of items that should actually be manufactured as part of the job; contains \texttt{<item>} elements = objects to be manufactured
    \end{itemize}
\end{itemize}